% v005 (2026-11-17): P3 and P3Net as implemented in v005 (p3net/v005, cards p3net.md, p3-canonical.md,
% p3net-surrogate.md). Every sentence about behaviour maps to a row of notes/archive/v005/plans/v005-claims.md:
% engine C-01..C-07; canonical C-08, C-11; P3Net C-10, C-12, C-13, C-49, C-51, C-52; variants C-14, C-48 (Z-G1).
% The v003 description (ridge surrogate with cross-validation, kappa, geometric levels with HV-based growth) is
% obsolete and not reused.
\section{P3 and P3Net}
\label{sec:proposed-optimizer}

\paragraph{The P3 engine.} P3 keeps a pyramid of populations without target sizes~\citep{goldman2014parameterless}.
Each iteration adds one new random solution, improves it with a first-improvement hill climber and inserts it into the
lowest level if it is not yet in the pyramid. The solution is then mixed with every level of the pyramid, whether or not
it improved, and it is promoted to the next level only if it dominates its value before mixing and is not yet in the
pyramid. Mixing with a level uses that level's linkage clusters, built with P3's entropy-based distance: all
singletons are used, the root is not, and children of a zero-distance merge are dropped. Clusters are applied smallest
first, ties in random order, each with a donor from the level that differs from the solution on the cluster. With two
objectives, a hill-climbing move and a promotion require Pareto dominance, and a mixing step is kept under the
MO-GOMEA acceptance rule~\citep{luong2014mogomea}: the new solution dominates the current one, has the same
objectives, or is not dominated by the elitist archive of all evaluated solutions. Canonical P3 evaluates every mixing
step that changes the solution and uses no surrogate. P3 has no other parameters.

\paragraph{P3Net.} P3Net keeps this engine and the hill climber on real evaluations, and replaces the evaluations
of mixing steps by a surrogate: mixing with a level runs on predictions, and the solution is evaluated for real once,
after the level. The predicted objectives of a solution reached by a chain of surrogate-accepted steps are, in log
scale, the real objectives of the last really evaluated solution on the chain plus the predicted changes of all steps
since. The surrogate is first
used after the first hill climb has finished; until then mixing uses real evaluations. With the surrogate replaced by
an exact oracle, P3Net reproduces the trajectory of canonical P3, so the surrogate is the only difference between them.

\paragraph{The pairwise surrogate.} The surrogate predicts the difference of both objectives between two solutions,
in log scale, from features of the pair: which variables differ and between which values, and products of such
features for pairs of variables in the same cluster of the current linkage tree, so that the tree decides which
interactions the model can represent. It is antisymmetric by construction. One random forest per objective, with the
library's default hyperparameters and nothing tuned, is trained on all ordered pairs of evaluated solutions, or on
3\,000 of them drawn at random when there are more. Because the cost $f_2$ is predicted like $f_1$, P3Net never sees
the cost of a solution it has not evaluated.

\paragraph{Variants.} Two variants isolate how the surrogate is used and what the linkage tree adds. P3Net-gate uses
the surrogate as a gate instead of a replacement: a mixing step is evaluated for real only when its predicted
objectives would be accepted, and acceptance is then decided on the real values. P3Net-nolink uses the same pairwise
model without the interaction features of the linkage tree.
